\begin{table}[htb]
\centering
\begin{tabular}{| c | D{.}{.}{-1} | D{.}{.}{-1} |}
\hline
%  {\bf Source of uncertainty}    & \multicolumn{1}{c|}{{\bf ${\delta}N_{SK}$/$N_{SK}$}} & \multicolumn{1}{c|}{{\bf ${\delta}N_{SK}$/$N_{SK}$}}\\
  {\bf Source of uncertainty}    & \multicolumn{1}{c|}{{\bf ${\delta}N_{SK}$/$ <N_{SK} >$}} & \multicolumn{1}{c|}{{\bf ${\delta}N_{SK}$/$<N_{SK}>$}}\\
                                 & \multicolumn{1}{c|}{($\beta=0$)} & \multicolumn{1}{c|}{($\beta=1$)}\\
\hline
  SuperK detector    				& 5.05\%                         & 3.03\%                        \\
  FSI + SI + PN 					& 4.34\%				& 2.45\%			\\
  SuperK detector + FSI + SI + PN     & 6.66\%                         & 3.90\%                        \\
\hline
  BANFF (prefit)                 &  12.37\%                        & 12.91\%                       \\
  BANFF (postfit)                & 9.72\%                         & 10.29\%                        \\
\hline
  BANFF (Non-oxygen, postfit)    	& 3.40\%         	         & 3.02\%                        \\
  BANFF (Oxygen-specific postfit) 	& 9.24\%                    & 10.14\%                        \\
  MEC (Oxygen-specific postfit)	   	& 8.86\%                        	 & 9.28\%                        \\
\hline
  Oscillation parameters         & 4.72\%                         & 13.38\%                       \\
\hline
  Total (BANFF prefit)           &  14.96\%                        & 18.60\%                       \\
  Total (BANFF postfit)          & 12.75\%                         & 17.21\%                       \\
\hline
\end{tabular}
\caption{
 Effect of 1$\sigma$ variations of the groups of systematics parameters on total number of
 1-ring electron-like events in the reverse-horn current sample.
 The BANFF parameters considered to be oxygen-specific are Oxygen MEC normalisation, Fermi Energy, Binding Energy, CC Coherent Normalisation,
 NC Coherent normalization and the nue/numu normalisation. 
The effect of each group of systematics is computed as the standard deviation divided by the mean of the distribution
of the total number of superK events.
\NumToys are performed.
% \newline
%{\bf TO UPDATE}
%\newline
}
\label{tab:nsk_systematic_table_summary}
\end{table}
